\chapter{Aplicacions, dominis, codominis i composició}\label{cap:exercicis-aplicacions}

Aquests exercicis proposats corresponen al capítol homònim de la part teòrica (capítol~\ref{cap:aplicacions} de la Part~\ref{part:teoria}) i es resolen amb les tècniques dels exercicis resolts de la part pràctica (capítol~\ref{cap:practica-aplicacions} de la Part~\ref{part:practica}). Els marcats amb una indicació són els més exigents.

\section{Llenguatge formal, composició i estabilitat}

\begin{enumerate}
  \item Escriviu amb els predicats \(\Dom\), \(\Cod\) i \(\Com\) que \(h\) és la composició de \(f\) seguida de \(g\), que el domini de \(h\) coincideix amb el de \(f\) i que el codomini de \(h\) coincideix amb el de \(g\).

  \item Traduïu al llenguatge dels predicats la frase \enquote{\(f\) i \(g\) tenen el mateix codomini}. Feu-ho de dues maneres: amb una variable auxiliar per al codomini comú i sense.

  \item Siguin \(f\colon A\to B\), \(g\colon B\to C\) i \(h\colon C\to D\). Escriviu les dues maneres d'agrupar \(h\circ g\circ f\) i justifiqueu-ne la igualtat a partir de l'axioma d'associativitat.

  \item Demostreu que si \(\Com(f,g,h)\) i \(\Com(f,g,h')\), aleshores \(h=h'\). \begin{indicacio}És la unicitat de la composició (proposició~\ref{prop:composicio-unica}); cap raonament nou, només citar-la correctament.\end{indicacio}
\end{enumerate}

\section{Identitats i conjunts}

\begin{enumerate}[resume]
  \item Demostreu que per a tot conjunt \(A\) es té \(\id_A\circ\id_A=\id_A\), i que \(\id_A\) és l'única aplicació \(A\to A\) amb aquesta propietat que a més és neutra per composició.

  \item Sigui \(e\colon A\to A\) una aplicació amb \(e\circ e=e\) (un \emph{idempotent}). Doneu, en el model clàssic, un exemple d'un idempotent que no sigui la identitat. \begin{indicacio}Una aplicació constant \(A\to A\) amb valor un punt fix serveix; comproveu \(e\circ e=e\).\end{indicacio}
\end{enumerate}

\section{Igualtat d'aplicacions i elements generalitzats}

\begin{enumerate}[resume]
  \item Siguin \(f,g\colon A\to B\). Demostreu que si \(f=g\), aleshores \(f\circ x=g\circ x\) per a tot element generalitzat \(x\colon X\to A\). (És la direcció fàcil del criteri d'igualtat.)

  \item Expliqueu amb un exemple del model clàssic per què la separació per elements \emph{globals} no seria suficient si no la postuléssim: descriviu dos conjunts i dues aplicacions paral·leles que un observador que només mirés punts de \(\mathbf1\) no podria distingir sense l'axioma. \begin{indicacio}Penseu en què passa quan el domini no té elements globals, com \(\mathbf0\); relacioneu-ho amb l'observació~\ref{obs:separacio-no-degeneracio}.\end{indicacio}
\end{enumerate}

\section{Isomorfismes}

\begin{enumerate}[resume]
  \item Sigui \(f\colon A\to B\) un isomorfisme. Demostreu que \({(f^{-1})}^{-1}=f\).

  \item Demostreu que la isomorfia es comporta transitivament: de \(A\cong B\) i \(B\cong C\), deduïu \(A\cong C\). Indiqueu quin isomorfisme concret exhibeix \(A\cong C\).

  \item Demostreu que \(\cong\) és una relació d'equivalència sobre els conjunts (reflexiva, simètrica i transitiva), tot dient quins resultats del capítol useu a cada propietat.

  \item Siguin \(f,g\colon A\to B\) isomorfismes amb \(f^{-1}=g^{-1}\). Demostreu que \(f=g\). \begin{indicacio}Componeu la igualtat de les inverses convenientment amb \(f\) o \(g\) i useu la neutralitat.\end{indicacio}
\end{enumerate}

\section{Terminals, inicials i no-degeneració}

\begin{enumerate}[resume]
  \item Demostreu que si \(I_1\) i \(I_2\) són tots dos inicials, hi ha un únic isomorfisme \(I_1\to I_2\). (Dual de la unicitat del terminal.)

  \item Demostreu que tota aplicació \(A\to\mathbf1\) és \({!}_A\). Deduïu-ne que, per a cada estadi \(X\), el conjunt \(\mathbf1\) té exactament un element generalitzat en l'estadi \(X\), i compareu-ho amb el nombre d'aplicacions \(\mathbf1\to A\) per a un conjunt \(A\) qualsevol. \begin{indicacio}El que és únic és l'aplicació \emph{cap} a \(\mathbf1\); les aplicacions \emph{des de} \(\mathbf1\) són els punts d'\(A\), i poden ser molts o cap.\end{indicacio}

  \item Demostreu que no existeix cap aplicació \(\mathbf1\to\mathbf0\) reproduint l'argument del resolt corresponent, i deduïu que \(\mathbf0\) no té cap element global.

  \item Suposeu un univers alternatiu en què se satisfan tots els axiomes tret de la no-degeneració, i en què \(\mathbf0\cong\mathbf1\). Demostreu que aleshores entre cada parell de conjunts hi ha exactament una aplicació. \begin{indicacio}Vegeu l'observació~\ref{obs:zero-u-sense-separacio}; useu que tot conjunt rebria i emetria una única fletxa cap a i des de l'objecte col·lapsat.\end{indicacio}

  \item Escriviu l'enunciat dual de la proposició~\ref{prop:no-fletxa-u-zero} i digueu si és cert. \begin{indicacio}El dual és la mateixa proposició: la no-degeneració és autodual. Compareu-ho amb l'exercici resolt sobre la no-autodualitat de la separació global.\end{indicacio}

  \item \textbf{Axiomes usats.} Per a cadascun dels resultats següents, digueu quins dels nou axiomes del capítol intervenen a la demostració i quins no hi intervenen: la unicitat de la composició (proposició~\ref{prop:composicio-unica}), la unicitat de la inversa (proposició~\ref{prop:inversa-unica}), la unicitat del terminal llevat d'isomorfisme (proposició~\ref{prop:terminal-unic}), la inexistència d'aplicacions \(\mathbf1\to\mathbf0\) (proposició~\ref{prop:no-fletxa-u-zero}) i l'extensionalitat global (proposició~\ref{prop:extensionalitat-global}). \begin{indicacio}Les tres primeres només usen els axiomes estructurals i, si escau, les definicions; les dues últimes són les úniques que necessiten la no-degeneració i la separació global, respectivament.\end{indicacio}
\end{enumerate}

\section{Elements globals i aplicacions constants}

\begin{enumerate}[resume]
  \item Sigui \(b\colon\mathbf1\to B\) un punt i \(c\colon A\to B\) la constant de valor \(b\). Demostreu que si \(A\) té almenys un element global, aleshores \(c\) queda determinada pel seu valor sobre qualsevol d'ells.

  \item Siguin \(b,b'\colon\mathbf1\to B\) dos punts i \(c,c'\colon A\to B\) les constants respectives. Demostreu que \(c=c'\) si i només si \(b=b'\), suposant que \(A\) no és inicial. \begin{indicacio}Per a una direcció componeu amb un element global de \(A\); l'altra és immediata. On s'usa que \(A\) no és inicial?\end{indicacio}
\end{enumerate}
